\section{可靠性、Plan Divergence 与 LLM OS}

课程最后把 agent 的剩余问题集中到 reliability。传统软件可以在相同输入下重复执行，agent policy 则会跨许多 stochastic step 与动态环境互动。局部成功率看似很高，长轨迹仍会快速累积失败。

\subsection{自主 Agent 的关键问题}

Key Issues 页列出 reliability、looping、plan divergence、benchmark、observability、security 与 human fallback。它们不是部署后的附加功能，而是 agent architecture 的一部分。

\lecturefigure{slide-109.jpg}{Autonomous agent 的可靠性、循环、测试与安全问题}{官方 CS25 V4 slide deck，第 109 页}

\begin{importantbox}{长轨迹成功率的乘法效应}
若每一步独立成功概率为 $p$，长度为 $T$ 的 trajectory 全部成功概率近似
\[
P(\text{all success})=p^T.
\]
即使 $p=0.99$，执行 100 步也只有约 $0.99^{100}\approx 36.6\%$ 的全程无错概率。真实错误并不独立，早期偏离还会改变后续 observation，使失败更严重。
\end{importantbox}

\teachervoice{讲者用“95\% 的 agent 仍有 5\% 会失败”说明软件部署差异：传统程序可以修复 deterministic bug，stochastic agent 的少数失败会在大规模用户和长轨迹中不断出现。}

\begin{knowledgebox}{Observability 要记录什么}
\begin{enumerate}
  \item 每一步 observation、selected action 与 tool response；
  \item model/prompt/tool version 和 memory snapshot；
  \item permission decision、user confirmation 与 side effect；
  \item verifier verdict、retry、fallback 与 final outcome。
\end{enumerate}
日志的目标是支持复现和定位，不是无限保存敏感内容；必须配合最小化、加密与 retention policy。
\end{knowledgebox}

\subsection{Plan divergence：偏离后怎样回到轨道}

Plan divergence 页用 ideal path 与 actual path 对比。Agent 在某一步误读页面、工具报错或目标分解错误后，后续 action 可能建立在错误 state 上，最终进入 loop。本节承接前一页的 reliability，进一步解释为什么系统必须周期性比较计划 state 与真实 observation。

\lecturefigure{slide-110.jpg}{Plan divergence：实际轨迹逐步偏离理想路径}{官方 CS25 V4 slide deck，第 110 页}

\begin{importantbox}{Recovery policy}
Agent 不应只问“下一步是什么”，还要周期性检查
\[
d_t=D(s_t,\hat{s}_t,g),
\]
其中 $s_t$ 是真实 observation，$\hat{s}_t$ 是计划预期 state，$g$ 是目标。当 divergence score $d_t$ 超过阈值时，系统应暂停、重新观察、回滚到 checkpoint、重规划或请求人工介入。
\end{importantbox}

\begin{warningbox}{Loop detector 不能只统计重复文本}
Agent 可能用不同措辞重复同一失败 action，也可能在页面间循环。检测应结合 action signature、state hash、error pattern、progress metric 与预算消耗。达到步数、时间或成本上限时必须硬停止。
\end{warningbox}

\teachervoice{讲者评价 AutoGPT 是很好的 prototype，却“实际上做不了什么”，因为它容易循环、偏离并持续犯错。这里的重点不是否定原型，而是指出缺少 error correction 的长轨迹系统无法靠继续生成自行恢复。}

\subsection{LLM OS：一个有用但不完整的系统类比}

Karpathy 的 LLM OS 图把 LLM 视为 CPU，把 context window 视为 RAM，把 embedding store 视为 file system，把 calculator、Python、terminal 视为传统软件或 ALU，把图像、音频、浏览器和其他模型视为 peripheral。

\lecturefigure{slide-111.jpg}{Karpathy 的 LLM OS 组件映射}{官方 CS25 V4 slide deck，第 111 页}

\begin{knowledgebox}{读图：LLM OS 映射表}
\begin{tabularx}{\textwidth}{L{0.22\textwidth} L{0.24\textwidth} X}
\toprule
传统系统 & LLM system 类比 & 尚缺的保证 \\
\midrule
CPU & Foundation model & 确定执行、类型化指令与正确性。 \\
RAM & Context window & 隔离、分页一致性与受控写入。 \\
File system & Retrieval/memory store & 事务、权限、版本与删除语义。 \\
ALU/software & Calculator、Python、tools & 沙箱、资源上限与错误语义。 \\
Peripheral/network & Vision、audio、browser、other agents & 认证、信任边界与输入验证。 \\
\bottomrule
\end{tabularx}
\end{knowledgebox}

\begin{warningbox}{“OS”这个词不会自动带来操作系统能力}
真正 OS 负责 scheduling、memory protection、process isolation、permission、interrupt 与 recovery。LLM orchestration framework 若没有这些机制，只是形状像 OS，而没有 OS 的安全和可靠性保证。
\end{warningbox}

\subsection{Generalized AI system 与 action engine}

下一页把用户 chat interface、action engine、agents、memory、tools 与外部世界连接成 generalized system。Action engine 负责 route、plan、permission 与 aggregation，是模型之外最关键的 control plane。

\lecturefigure{slide-112.jpg}{Chat interface、action engine、agents、tools 与 memory 的系统图}{官方 CS25 V4 slide deck，第 112 页}

\teachervoice{结尾愿景把 chat interface 后面的 action engine 设为路由中心：它把任务分给不同 agent，再汇总结果。这个位置也应承担权限、预算、验证和中止，而不能只是一个 prompt router。}

\begin{importantbox}{Control plane 与 data plane}
\begin{itemize}
  \item \textbf{Control plane}：解析目标、选择 agent/tool、分配权限、维护 state、决定 retry/fallback；
  \item \textbf{Data plane}：实际执行 search、code、browser、database、message 或 device action。
\end{itemize}
把两者分离后，可以限制某个 model proposal 不直接等于真实 side effect。
\end{importantbox}

\subsection{Future needs：error correction、permission 与 sandbox}

最后一页没有给出“agent 已解决”的结论，而是列出未完成的基础设施：error correction、better framework、security、user permission、risky-domain stability 与 sandboxing。读图时应把这些项目视为 deployment precondition，并与前面的 long-horizon failure、权限和 side effect 逐项对应。

\lecturefigure{slide-113.jpg}{AI agent 仍需 error correction、security、permission 与 sandbox}{官方 CS25 V4 slide deck，第 113 页}

\begin{importantbox}{高风险 action 的最小安全栈}
\begin{enumerate}
  \item \textbf{Least privilege}：只授予完成当前任务所需的最小 credential 和 action scope；
  \item \textbf{Confirmation gate}：支付、发送、删除、发布和权限变更需要显式确认；
  \item \textbf{Sandbox}：代码、文件和网络访问在隔离环境中运行；
  \item \textbf{Budget}：限制 token、步骤、时间、费用和外部调用；
  \item \textbf{Audit and recovery}：记录副作用，支持取消、补偿或人工接管。
\end{enumerate}
\end{importantbox}

\teachervoice{课程的最终落点不是“agent 很快会全自动”，而是 error correction、security、user permission 与 sandbox 仍未解决。尤其在 finance、legal 等高风险场景，稳定和安全必须先于更长 autonomy。}

\subsection{本章小结}

Agent reliability 受长轨迹乘法效应和 plan divergence 支配。LLM OS 提供了组织组件的语言，却不会自动产生 isolation、transaction 或 recovery。可部署系统需要 control plane、least privilege、verification、budget、fallback 与 sandbox 共同约束模型 proposal。

